% =====================================================================
%  Partie A — Questions à choix multiples
%
%  Modèle :
%    \begin{mcq}[points=1]
%      Énoncé (texte, code, image…)
%      \begin{choices}[columns=2]    % columns : facultatif
%        \wrong   ...
%        \correct ...                % bonne réponse (case noircie au corrigé)
%      \end{choices}
%    \end{mcq}
% =====================================================================

\exampart{Choix multiples}

\begin{instructions}
  Choisissez la bonne réponse et cochez-la dans la grille du cahier de
  réponses.
\end{instructions}

% --- Question sans code ---------------------------------------------
\begin{mcq}
  Quelle balise HTML crée un lien hypertexte vers une autre page ?
  \begin{choices}[columns=4]
    \wrong   \codeinline{<link>}
    \correct \codeinline{<a>}
    \wrong   \codeinline{<href>}
    \wrong   \codeinline{<nav>}
  \end{choices}
\end{mcq}

% --- Python ---------------------------------------------------------
\begin{mcq}
  Qu'affiche le programme Python suivant ?
\begin{code}{python}
def somme_pairs(valeurs):
    """Additionne les nombres pairs de la liste."""
    total = 0
    for v in valeurs:
        if v % 2 == 0:
            total += v
    return total

print(somme_pairs([1, 2, 3, 4, 5, 6]))  # Résultat affiché ?
\end{code}
  \begin{choices}[columns=4]
    \wrong   \codeinline{21}
    \correct \codeinline{12}
    \wrong   \codeinline{9}
    \wrong   \codeinline{6}
  \end{choices}
\end{mcq}

% --- JavaScript -----------------------------------------------------
\begin{mcq}
  Que contient la constante \codeinline{resultat} après l'exécution de ce code ?
\begin{code}{javascript}
const nombres = [3, 1, 4, 1, 5];
const resultat = nombres
  .filter(n => n > 2)
  .map(n => n * 2);
console.log(resultat);
\end{code}
  \begin{choices}[columns=2]
    \wrong   \codeinline{[6, 2, 8, 2, 10]}
    \correct \codeinline{[6, 8, 10]}
    \wrong   \codeinline{[3, 4, 5]}
    \wrong   \codeinline{[8, 10]}
  \end{choices}
\end{mcq}

% --- HTML + CSS -----------------------------------------------------
\begin{mcq}
  Soit le fragment HTML et la feuille de style ci-dessous. De quelle
  couleur le texte « Attention ! » est-il affiché ?
\begin{code}{html}
<p id="message" class="alerte">Attention !</p>
\end{code}
\begin{code}{css}
/* Feuille de style */
p        { color: black; }
.alerte  { color: orange; }
#message { color: red; }
p.alerte { color: blue; }
\end{code}
  \begin{choices}[columns=4]
    \wrong   Noir
    \wrong   Orange
    \correct Rouge
    \wrong   Bleu
  \end{choices}
  % Le sélecteur d'identifiant #message a la plus grande spécificité.
\end{mcq}

% --- Node.js ---------------------------------------------------------
\begin{mcq}
  Dans quel ordre les lettres sont-elles affichées par ce programme
  Node.js ? (Le fichier \codeinline{donnees.txt} existe.)
\begin{code}{nodejs}
const fs = require('fs');

console.log('A');
fs.readFile('donnees.txt', 'utf8', (err, data) => {
  console.log('B');
});
console.log('C');
\end{code}
  \begin{choices}[columns=4]
    \wrong   A, B, C
    \correct A, C, B
    \wrong   B, A, C
    \wrong   C, A, B
  \end{choices}
  % readFile est asynchrone : son rappel s'exécute après le code synchrone.
\end{mcq}

% --- Java : choix sous forme de code ---------------------------------
\begin{mcq}
  Laquelle de ces déclarations Java compile sans erreur ?
  \begin{choices}
    \wrong   \codeinline{int x = "5";}
    \wrong   \codeinline{String s = 'abc';}
    \correct \codeinline{double d = 3;}
    \wrong   \codeinline{boolean b = 1;}
  \end{choices}
\end{mcq}
